\subsubsection{Pseudocódigo.}

\begin{tcolorbox}[colback=white, colframe=azul_unam, title=Pseudocódigo: sinc\_lamport.c (P2P), breakable]

\textbf{Algoritmo} Sincronizacion\_Lamport

% ===================== DEFINICIONES Y ESTRUCTURAS =====================

\textbf{Constantes y Estructuras (lamport\_p2p.h)}

\begin{algorithmic}[1]

\State \textbf{Definir} $NUM\_NODOS \gets 3$
\State \textbf{Definir} $NUM\_EVENT0 \gets 7$
\State \textbf{Definir} $CMD\_EX \gets$ "EXIT"
\State \textbf{Definir} $CMD\_EN \gets$ "ENVIAR"
\State \textbf{Definir} $CMD\_IN \gets$ "LOCAL"
\State \textbf{Definir} $DIRECTION \gets$ "127.0.0.1"

\vspace{0.2cm}

\State \textbf{Estructura} $Mensaje\_L$
\State \quad $id\_proceso \in \mathbb{Z}$
\State \quad $indice \in \mathbb{Z}$
\State \quad $reloj \in \mathbb{Z}$
\State \quad $peticion \in \text{Alfabeto}[200]$
\State \textbf{FinEstructura}

\end{algorithmic}

\vspace{0.3cm}

% ===================== VARIABLES GLOBALES =====================

\textbf{Variables Globales}

\begin{algorithmic}[1]
\State \textbf{Definir} $instancia \in Mensaje\_L$
\State \textbf{Definir} $directorio \in \mathbb{Z}[NUM\_NODOS]$
\end{algorithmic}

\vspace{0.3cm}

% ===================== MAIN =====================
\begin{algorithmic}[1]
    
\Procedure{main}{}

\State $pid \gets fork()$

\If{$pid < 0$}
    \State \textbf{Salir}(1)
\EndIf

% ===================== PROCESO HIJO / NIETOS =====================
\If{$pid = 0$} \Comment{Proceso hijo creador de nodos P2P}

    \For{$i \gets 0$ \textbf{hasta} $NUM\_NODOS - 1$}
        \State $directorio[i] \gets 3000 + i$
    \EndFor

    \For{$i \gets 0$ \textbf{hasta} $NUM\_NODOS - 1$}
        
        \If{$fork() = 0$} \Comment{Proceso nieto (Nodo P2P individual)}
            
            \State $instancia.id\_proceso \gets 5 * (i + 1)$
            \State $instancia.indice \gets i$
            \State $instancia.reloj \gets 0$
            \State LimpiarMemoria$(instancia.peticion)$
            
            \State conexion\_p2p$(instancia, directorio)$
            
            \State \textbf{Salir}(0)
            
        \EndIf
        
    \EndFor

    \For{$i \gets 0$ \textbf{hasta} $NUM\_NODOS - 1$}
        \State EsperarHijo$()$
    \EndFor

    \State \textbf{Salir}(0)

% ===================== PROCESO PADRE =====================
\Else \Comment{Proceso padre inyector de comandos}

    \State GenerarSemilla$(TiempoActual() \text{ XOR } (ObtenerPID() \ll 16))$
    \State DormirSegundos$(1)$
    
    \State Escribir("=========================================")
    \State Escribir(" [PADRE] INICIANDO SIMULACIÓN P2P (RELOJES DE LAMPORT) ")
    \State Escribir("=========================================")
    
    \State \textbf{Definir} $inyeccion \in Mensaje\_L$
    
    \For{$i \gets 0$ \textbf{hasta} $NUM\_EVENT0 - 1$}
        
        \State $indice\_elegido \gets Aleatorio() \text{ MOD } NUM\_NODOS$
        \State $puerto\_elegido \gets 3000 + indice\_elegido$
        
        \State LimpiarMemoria$(inyeccion)$
        \State $inyeccion.id\_proceso \gets 99$
        \State $inyeccion.reloj \gets 0$ \Comment{El proceso padre no altera el tiempo}
        
        \If{$(Aleatorio() \text{ MOD } 2) = 1$}
            \State CopiarCadena$(inyeccion.peticion, CMD\_EN)$
        \Else
            \State CopiarCadena$(inyeccion.peticion, CMD\_IN)$
        \EndIf
        
        \State Escribir("\textbackslash n[PADRE] Ordenando a la computadora ", $(indice\_elegido + 1) * 5$, " con el puerto ", $puerto\_elegido$, " a realizar un evento...")
        
        \State enviar\_msj$(puerto\_elegido, inyeccion)$
        \State DormirMicrosegundos$(100000)$
        
    \EndFor

    \State Escribir("\textbackslash n[PADRE] Simulación terminada. Esperando a que la red se estabilice...")
    
    \State Escribir("\textbackslash n======================================")
    \State Escribir(" [PADRE] INICIANDO PROTOCOLO DE APAGADO GLOBAL ")
    \State Escribir("========================================\textbackslash n")

    \For{$i \gets 0$ \textbf{hasta} $NUM\_NODOS - 1$}
        
        \State LimpiarMemoria$(inyeccion)$
        \State $inyeccion.id\_proceso \gets 99$
        \State $inyeccion.reloj \gets 0$
        \State CopiarCadena$(inyeccion.peticion, CMD\_EX)$
        
        \State Escribir("\textbackslash n[PADRE] Enviando comando de apagado a puerto ", $3000 + i$, "...")
        \State enviar\_msj$(3000 + i, inyeccion)$
        \State DormirMicrosegundos$(50000)$
        
    \EndFor

    \State EsperarHijo$()$
    \State Escribir("\textbackslash n[SISTEMA GLOBAL] Todos los nodos P2P han finalizado exitosamente. Fin del programa.")

\EndIf

\State \Return 0

\EndProcedure

\end{algorithmic}
\end{tcolorbox}

\begin{tcolorbox}[colback=white, colframe=azul_unam, title=Pseudocódigo: cnx\_lamport.c (Conexión P2P), breakable]

\textbf{Algoritmo} Conexion\_P2P

% ===================== VARIABLES GLOBALES =====================

\textbf{Constantes y Variables Globales}

\begin{algorithmic}[1]

\State \textbf{Definir} $PORT \in \mathbb{N}, PORT \gets 0$
\State \textbf{Definir} $local, externo \in Mensaje\_L$
\State \textbf{Definir} $puertos \in \mathbb{Z}[NUM\_NODOS]$
\State \textbf{Definir} $encendido \gets 1$
\State \textbf{Definir} $enviar \gets 0$

\vspace{0.2cm}

% ===================== FUNCIÓN MAYOR =====================

\Procedure{mayor}{a, b}
    \If{$a > b$}
        \State \Return $a$
    \Else
        \State \Return $b$
    \EndIf
\EndProcedure

\vspace{0.3cm}

% ===================== CONEXION P2P =====================

\Procedure{conexion\_p2p}{user, dir}

    \State \textbf{Definir} $thread\_id \in \text{Hilo}$
    
    \State CopiarMemoria$(puertos, dir, \text{sizeof}(puertos))$
    \State $PORT \gets dir[user.indice]$
    \State $local.id\_proceso \gets user.id\_proceso$
    \State $local.indice \gets user.indice$
    \State $local.reloj \gets local.reloj + 1$ \Comment{Inicialización del reloj local}

    \State CrearHilo$(thread\_id, \text{hilo\_escucha}, NULL)$
    
    \State GenerarSemilla$(\text{Tiempo\_Actual}() \text{ XOR } ObtenerPID())$

    \While{$encendido$}
    
        \If{$enviar = 1$}
            \State $enviar \gets 0$
            
            \State $local.reloj \gets local.reloj + 1$ \Comment{REGLA LAMPORT: Avanzar reloj antes de enviar}
            
            \State \textbf{Definir} $indice\_new \in \mathbb{Z}$
            \Repeat
                \State $indice\_new \gets Aleatorio() \text{ MOD } NUM\_NODOS$
            \Until{$indice\_new \neq local.indice$}
            
            \State Escribir("\textbackslash n[C", $local.id\_proceso$, "] Enviare mensaje a la computadora con puerto ", $puertos[indice\_new]$, "...")
            
            \State LimpiarMemoria$(local.peticion)$
            \State FormatearCadena$(local.peticion, \text{"Soy la computadora \%02d"}, local.id\_proceso)$
            
            \State enviar\_msj$(puertos[indice\_new], local)$
        \EndIf
        
    \EndWhile

    \State EsperarHilo$(thread\_id)$

\EndProcedure

\vspace{0.3cm}

% ===================== HILO ESCUCHA (SERVIDOR TCP) =====================

\Procedure{hilo\_escucha}{arg}

    \State \textbf{Definir} $server\_fd, new\_socket \in \mathbb{Z}$
    \State \textbf{Definir} $address \in sockaddr\_in$
    \State \textbf{Definir} $opt \gets 1$

    \State $server\_fd \gets \text{socket}(AF\_INET, SOCK\_STREAM, 0)$
    \State \text{setsockopt}(server\_fd, SOL\_SOCKET, SO\_REUSEADDR, opt)

    \State $address.sin\_family \gets AF\_INET$
    \State $address.sin\_addr.s\_addr \gets INADDR\_ANY$
    \State $address.sin\_port \gets htons(PORT)$

    \State \text{bind}(server\_fd, address)
    \State \text{listen}(server\_fd, 3)

    \State Escribir("[C", $local.id\_proceso$, "] Escuchando Mensaje\_L en puerto ", $PORT$, "...")

    \While{$encendido$}
        
        \State $addrlen \gets \text{sizeof}(address)$
        \State $new\_socket \gets \text{accept}(server\_fd, address, addrlen)$

        \If{$new\_socket < 0$}
            \State \textbf{Romper}
        \EndIf

        \State LimpiarMemoria$(externo)$
        \State Leer$(new\_socket, externo)$

        \If{$externo.id\_proceso = 99$} \Comment{Mensajes del Padre (Control)}
            
            \If{CompararCadena$(externo.peticion, CMD\_EX) = 0$}
                \State Escribir("\textbackslash n[C", $local.id\_proceso$, "] Apagando por comando remoto...")
                \State $encendido \gets 0$
            \ElsIf{CompararCadena$(externo.peticion, CMD\_EN) = 0$}
                \State $enviar \gets 1$
            \ElsIf{CompararCadena$(externo.peticion, CMD\_IN) = 0$}
                \State Escribir("\textbackslash n[C", $local.id\_proceso$, "] -> EVENTO INTERNO <- Ejecutando tarea local...")
                \State $local.reloj \gets local.reloj + 1$
            \EndIf
            
        \Else \Comment{Mensajes de otros Nodos P2P}
            \State $local.reloj \gets mayor(local.reloj, externo.reloj) + 1$ \Comment{REGLA LAMPORT: Al recibir}
        \EndIf

        \State Escribir("\textbackslash nInicio -> Mensaje Computadora [C", $local.id\_proceso$, "] <-")
        \State Escribir("--- ID\_proceso: ", $externo.id\_proceso$)
        \State Escribir("--- Peticion: ", $externo.peticion$)
        \State Escribir("--- Reloj Externo (Ts): ", $externo.reloj$)
        \State Escribir("--- Reloj Local Actualizado: ", $local.reloj$)
        \State Escribir("Fin -> Mensaje Computadora [C", $local.id\_proceso$, "] <-\textbackslash n")

        \State Cerrar$(new\_socket)$

    \EndWhile

    \State Cerrar$(server\_fd)$
    \State \Return NULL

\EndProcedure

\vspace{0.3cm}

% ===================== ENVIAR MENSAJE (CLIENTE TCP) =====================

\Procedure{enviar\_msj}{puerto\_destino, datos}

    \State \textbf{Definir} $sock \in \mathbb{Z}, sock \gets 0$
    \State \textbf{Definir} $serv\_addr \in sockaddr\_in$

    \If{$(sock \gets \text{socket}(AF\_INET, SOCK\_STREAM, 0)) < 0$}
        \State \Return
    \EndIf

    \State $serv\_addr.sin\_family \gets AF\_INET$
    \State $serv\_addr.sin\_port \gets htons(puerto\_destino)$
    \State inet\_pton$(AF\_INET, DIRECTION, serv\_addr.sin\_addr)$

    \If{\text{connect}(sock, serv\_addr) < 0}
        \State Escribir("Error de conexion")
        \State Cerrar$(sock)$
        \State \Return
    \EndIf

    \State Enviar$(sock, datos)$
    \State Cerrar$(sock)$

\EndProcedure

\end{algorithmic}
\end{tcolorbox}